\section*{\S\,21. Der Fundamentalsatz der algebraischen Erweiterungen.}

\textbf{Definition 1. Hauptfunktionen und Hauptfaktoren der regulären Wohlordnung.} -- Es sei $\mathfrak N=\K_\sigma$ eine algebraische Erweiterung des wohlgeordneten Körpers $\K_0$ und in bezug auf $\K_0$ regulär geordnet, $\mathfrak J=\mathfrak J_\sigma=\{j_\nu\}$ $(0\leq\nu<\sigma)$ das System der Hauptelemente, $\K_\nu$ bzw. $\mathfrak J_\nu$ der durch $j_\nu$ bestimmte Hauptabschnitt von $\mathfrak N$ bzw. Abschnitt von $\mathfrak J$. Das Element $j_\nu$ ist Nullstelle einer bestimmten irreduziblen Funktion $f_\nu(x)$ aus $\K_0$ (mit $\varepsilon$ als Koeffizient der höchsten Potenz), ebenso Nullstelle einer bestimmten irreduziblen Funktion $\varphi_\nu(x)$ des Körpers $\K_\nu$, welche gleich $f_\nu(x)$ sein kann und für $\nu=0$ stets gleich $f_\nu(x)$, in jedem Falle aber ein Faktor von $f_\nu(x)$ ist. Wir gelangen so zu zwei Funktionensystemen
\[
F_\sigma=F=\{f_\nu(x)\},\qquad \Phi_\sigma=\Phi=\{\varphi_\nu(x)\}\qquad (0\leq \nu<\sigma).
\]
Wir nennen die Funktionen $f_\nu(x)$ \emph{Hauptfunktionen}, die Funktionen $\varphi_\nu(x)$ \emph{Hauptfaktoren} der regulären Wohlordnung und bezeichnen den durch $f_\nu(x)$ bzw. $\varphi_\nu(x)$ ($\nu\geq1$) bestimmten Abschnitt von $F$ bzw. $\Phi$ mit $F_\nu$ bzw. $\Phi_\nu$.\footnote{Um Mißverständnisse zu vermeiden, weisen wir noch darauf hin, daß den Funktionen $f_\nu(x)$ innerhalb $F$ eine bestimmte (durch den Index bezeichnete) Reihenfolge zukommt, die natürlich nichts mit der Reihenfolge zu tun hat, welche den Funktionen auf Grund von \S\,20, Def. 4 in $\K_0$ erteilt wurde.} -- Unter Bezugnahme auf diese Bezeichnung ergibt sich nun leicht:

\textbf{1.} \emph{Sind sämtliche Hauptfunktionen normal (in bezug auf $\K_0$), so ist der regulär geordnete Körper $\mathfrak N=\K_\sigma$ selbst normal. Er ist gerade hinreichend groß, um das System $F$ der Hauptfunktionen vollständig zu zerfällen; ebenso ist jeder Hauptabschnitt $\K_\nu$ normal und reicht zur Zerfällung von $F_\nu$ gerade hin. Das System $F$ ist fortschreitend.}

Da nämlich jede der Normalfunktionen $f_\rho(x)$ $(0\leq \rho<\nu)$ aus $F_\nu$ $(0<\nu\leq\sigma)$ im Körper $\K_\nu=\K_0(\mathfrak J_\nu)$ eine Nullstelle $j_\rho$ besitzt, so muß sie in $\K_\nu$ vollständig zerfallen, und da der Körper $\K_\nu$ aus $\K_0$ nur durch Adjunktion von Nullstellen der Funktionen aus $F_\nu$ hervorgeht, so reicht er zur Zerfällung von $F_\nu$ gerade hin, ist also normal. Da ferner im Falle $\nu<\sigma$ $\K_\nu$ das Element $j_\nu$ nicht mehr enthält, die Funktion $f_\nu(x)$ in $\K_\nu$ also nicht vollständig zerfällt, so ist $F$ ein fortschreitendes System.

\textbf{Definition 2. Normale Wohlordnung.} -- Es bezeichne $\mathfrak N$ eine normale Erweiterung des wohlgeordneten Körpers $\K_0$; $\mathfrak N$ sei in bezug auf $\K_0$ regulär geordnet, die Hauptfunktionen seien sämtlich normal. Verwenden wir die Buchstaben in derselben Bedeutung wie vorher, so ist $j_\nu$ $(0\leq\nu<\sigma)$ Nullstelle der irreduziblen Funktion $\varphi_\nu(x)$ aus $\K_\nu$, und $\varphi_\nu(x)$ Faktor der in $\K_0$ irreduziblen Funktion $f_\nu(x)$. Da $\K_\nu$ ein wohlgeordneter Körper ist, so kommt auf Grund von Def. 4, \S\,20 auch den Funktionen innerhalb $\K_\nu$ eine bestimmte Anordnung zu. Dies gilt also auch von den irreduziblen Faktoren, welche $f_\nu(x)$ innerhalb $\K_\nu$ besitzt. Wir wollen nun annehmen, daß jede Funktion $\varphi_\nu(x)$, sofern sie nicht einziger irreduzibler Faktor von $f_\nu(x)$ innerhalb $\K_\nu$ ist, den übrigen irreduziblen Faktoren von $f_\nu(x)$ innerhalb $\K_\nu$ vorangeht. Ist diese Voraussetzung erfüllt, so sagen wir: $\mathfrak N$ ist \emph{in bezug auf $\K_0$ normal geordnet}. Die normale Wohlordnung ist also ein spezieller Fall der regulären.

\textbf{Definition 3. Ausgezeichnete Nullstellen der Hauptfunktionen, die zu einer normalen Wohlordnung gehören.} -- Es sei $\mathfrak N$ eine normale und normal geordnete Erweiterung des wohlgeordneten Körpers $\K_0$, und es mögen die bisher gebrauchten Zeichen ihre Bedeutung beibehalten. Da die Hauptfunktion $f_\nu(x)$ in bezug auf $\K_0$ normal ist, so kann der Körper $\K_\nu(j_\nu)$ erhalten werden, indem man $\K_\nu$ irgendeine Nullstelle von $f_\nu(x)$ adjungiert. Alle diese Nullstellen sind primitive Elemente des Körpers $\K_\nu(j_\nu)$ in bezug auf $\K_\nu$, ihr Grad ist gleich dem Grade des Hauptfaktors $\varphi_\nu(x)$, welcher jedenfalls größer als 1 ist. Im Falle unvollkommener Körper brauchen die Nullstellen von $\varphi_\nu(x)$ nicht alle verschieden zu sein. Nehmen wir also an, $\varphi_\nu(x)$ habe $r$ verschiedene Nullstellen $j_\nu,j_\nu',\ldots,j_\nu^{(r-1)}$ und sei vom Grade $n$, so ist $n\geq r\geq1$, $n>1$. Ist nun $i_\nu$ irgendeine Nullstelle von $f_\nu(x)$ (im Körper $\mathfrak N$), so erhalten wir sämtliche Elemente des Körpers $\K_\nu(i_\nu)$, indem wir in sämtlichen ganzen Funktionen $g(x)$ aus $\K_\nu$, deren Grad kleiner als $n$ ist, $x$ durch $i_\nu$ ersetzen. Unter diesen Funktionen gibt es also $r$, welche für $x=i_\nu$ in Nullstellen von $\varphi_\nu(x)$ übergehen. Diesen $r$ Funktionen kommt nach Def. 4, \S\,20 innerhalb $\K_\nu$ eine bestimmte Reihenfolge zu. Ist $\psi(x)$ die erste unter ihnen und ist $\psi(i_\nu)=j_\nu$, so nennen wir $i_\nu$ eine \emph{ausgezeichnete Nullstelle} von $f_\nu(x)$. (Diese Bezeichnung gebrauchen wir auch, wenn $r=1$ ist, was nur bei unvollkommenen Körpern vorkommen kann.) Es gilt der Satz:

\textbf{2.} \emph{Das Hauptelement $j_\nu$ ist ausgezeichnete Nullstelle von $f_\nu(x)$, und unter den Nullstellen des Hauptfaktors $\varphi_\nu(x)$ gibt es sonst keine ausgezeichnete.}

Da nämlich die Nullstellen $j_\nu,j_\nu',\ldots,j_\nu^{(r-1)}$ von $\varphi_\nu(x)$ in $\K_\nu$ nicht vorkommen, so sind die $r$ Funktionen aus $\K_\nu$, welche für $x=i_\nu$ in diese Nullstellen übergehen, alle wenigstens vom Grade 1. Aus dem Umstande, daß der wohlgeordnete Körper $\K_\nu$ mit den Elementen $0,\varepsilon$ beginnt, folgt, daß unter allen Funktionen aus $\K_\nu$, deren Grad $\geq1$ ist, $x$ die erste ist. Wählen wir also für $i_\nu$ eine der Nullstellen von $\varphi_\nu(x)$, so ist $x$ die erste Funktion aus $\K_\nu$, welche für $x=i_\nu$ in eine Nullstelle von $\varphi_\nu(x)$ übergeht. Diese ist aber dann und nur dann gleich $j_\nu$, wenn wir $i_\nu=j_\nu$ wählen, womit die Behauptung erwiesen ist. -- Aus 2. ergibt sich noch:

\textbf{3.} \emph{Im Falle $\varphi_\nu(x)=f_\nu(x)$, also jedenfalls für $\nu=0$, ist $j_\nu$ die einzige ausgezeichnete Nullstelle von $f_\nu(x)$.}

Nunmehr beweisen wir:

\textbf{4.} \emph{Ist $\mathfrak N=\K_0(\overline{\mathfrak J})$ normale Erweiterung des wohlgeordneten Körpers $\K_0$, ist jedes Element $i_\nu$ des Systems $\mathfrak J_\sigma=\overline{\mathfrak J}=\{i_\nu\}$ $(0\leq\nu<\sigma)$ Nullstelle einer Normalfunktion $f_\nu(x)$ aus $\K_0$ und ist das Funktionensystem $F_\sigma=F=\{f_\nu(x)\}$ ein fortschreitendes (oder, was hier dasselbe besagt: ist das System $\overline{\mathfrak J}$ ein fortschreitendes), so kann der Körper $\mathfrak N$ auf eine und nur eine Weise in bezug auf $\K_0$ derartig normal geordnet werden, daß $F$ das System der Hauptfunktionen darstellt und die Elemente $i_\nu$ ausgezeichnete Nullstellen der zugehörigen Hauptfunktionen $f_\nu(x)$ werden.}

Wir zeigen zunächst, daß es nicht zwei verschiedene Normalordnungen der bezeichneten Art geben kann. Nehmen wir an, es lägen zwei solche Normalordnungen $\mathfrak N'$ und $\mathfrak N''$ des Körpers $\mathfrak N$ vor. Wir bezeichnen mit $\mathfrak J_\sigma'=\mathfrak J'=\{j_\nu'\}$, $\{\K_\nu'\}$, $\Phi_\sigma'=\Phi'=\{\varphi_\nu'(x)\}$ die Systeme der Hauptelemente, Hauptabschnitte, Hauptfaktoren von $\mathfrak N'$, mit $\mathfrak J_\sigma''=\mathfrak J''=\{j_\nu''\}$, $\{\K_\nu''\}$, $\Phi_\sigma''=\Phi''=\{\varphi_\nu''(x)\}$ die entsprechenden Systeme für $\mathfrak N''$. Mit $\overline{\mathfrak J}_\nu$, $F_\nu$, $\mathfrak J_\nu'$, $\mathfrak J_\nu''$, $\Phi_\nu'$, $\Phi_\nu''$ seien die durch $i_\nu$, $f_\nu(x)$ u. s. f. bestimmten Abschnitte von $\overline{\mathfrak J}$, $F$ u. s. f. bezeichnet. Dann stellen $\K_\nu'$, $\K_\nu''$ zwei Normalordnungen desselben Körpers $\K_0(\overline{\mathfrak J}_\nu)=\K_0(\mathfrak J_\nu')=\K_0(\mathfrak J_\nu'')$ dar, welcher zur Zerfällung des Systems $F_\nu$ gerade ausreicht. Da $i_0$ ausgezeichnetes Element von $f_0(x)$ sein soll, so ist nach 3. $j_0'=i_0=j_0''$. Wenn die Gleichung $j_\nu'=j_\nu''$ nicht für jeden Index besteht, so nehmen wir den kleinsten, für welchen $j_\nu'\neq j_\nu''$ ist. Dann ist $\mathfrak J_\nu'=\mathfrak J_\nu''$, und da $\mathfrak J_\nu'$ und $\mathfrak J_\nu''$ die Systeme der Hauptelemente der in bezug auf $\K_0$ regulär geordneten Körper $\K_\nu'$, $\K_\nu''$ darstellen, so ist (auch unter Berücksichtigung der Anordnung) $\K_\nu'=\K_\nu''$ (\S\,20, 10.). Hieraus folgt weiter $\varphi_\nu'(x)=\varphi_\nu''(x)$, weil $\varphi_\nu'(x)$ bzw. $\varphi_\nu''(x)$ den ersten irreduziblen Faktor von $f_\nu(x)$ innerhalb $\K_\nu'=\K_\nu''$ darstellt. Ist ferner innerhalb $\K_\nu'=\K_\nu''$ $\psi(x)$ die erste Funktion, welche für $x=i_\nu$ Nullstelle von $\varphi_\nu'(x)=\varphi_\nu''(x)$ wird, so erhalten wir, weil $i_\nu$ ausgezeichnete Nullstelle von $f_\nu(x)$ sein soll, $j_\nu'=\psi(i_\nu)=j_\nu''$, im Widerspruch mit der Annahme $j_\nu'\neq j_\nu''$. Mithin ist $\mathfrak J'=\mathfrak J''$, und daraus folgt nach \S\,20, 10. $\mathfrak N'=\mathfrak N''$.

Wir haben jetzt zu zeigen, daß eine Normalordnung der verlangten Art existiert. -- Für $\sigma=1$ erhalten wir dieselbe, indem wir $\mathfrak N$ in bezug auf $\K_0$ primitiv mit $i_0$ als Hauptelement ordnen. Für $\sigma>1$ ist der Beweis durch Induktion zu führen. Ist $\sigma$ keine Limeszahl, so bezeichne $\K_{\sigma-1}$ diejenige Normalordnung des Körpers $\K_0(\overline{\mathfrak J}_{\sigma-1})$, für welche $F_{\sigma-1}$ das System der Hauptfunktionen darstellt und jedes Element $i_\nu$ $(\nu<\sigma-1)$ ausgezeichnete Nullstelle von $f_\nu(x)$ wird. Ist dann innerhalb des wohlgeordneten Körpers $\K_{\sigma-1}$ $\varphi_{\sigma-1}(x)$ der erste irreduzible Faktor von $f_{\sigma-1}(x)$, $\psi(x)$ die erste Funktion, welche für $x=i_{\sigma-1}$ zu einer Nullstelle von $\varphi_{\sigma-1}(x)$ wird, so setzen wir $\psi(i_{\sigma-1})=j_{\sigma-1}$ und ordnen $\mathfrak N$ primitiv in bezug auf $\K_{\sigma-1}$ mit $j_{\sigma-1}$ als Hauptelement. Dann hat man, wie sofort ersichtlich, die verlangte Normalordnung. -- Ist $\sigma$ Limeszahl, so stelle für jedes $\nu<\sigma$ $\K_\nu$ diejenige Normalordnung des Körpers $\K_0(\overline{\mathfrak J}_\nu)$ dar, für welche $F_\nu$ das System der Hauptfunktionen ist, während jedes Element $i_\mu$ $(\mu<\nu)$ ausgezeichnete Nullstelle von $f_\mu(x)$ wird. Dann zeigt eine einfache Überlegung, daß jeder Körper der Folge $\{\K_\nu\}$ $(0\leq\nu<\sigma)$ Abschnitt aller späteren ist, und daß die Gesamtheit der in diesen Körpern auftretenden Elemente eine Wohlordnung von $\mathfrak N$ liefert, die allen Bedingungen der Aufgabe entspricht.

\textbf{5.} \emph{Sind $\K_0$ und $\overline{\K}_0$ ähnlich-isomorphe Körper, ist $\mathfrak N=\K_\sigma$ eine normale und normal geordnete Erweiterung von $\K_0$, $\overline{\mathfrak N}=\overline{\K}_\sigma$ eine ebensolche von $\overline{\K}_0$, sind $F=\{f_\nu(x)\}$, $\overline F=\{\overline f_\nu(x)\}$ $(0\leq\nu<\sigma)$ die zugehörigen Systeme der Hauptfunktionen und werden durch die zwischen $\K_0$ und $\overline{\K}_0$ bestehende ähnlich-isomorphe Beziehung $\pi_0$ die Funktionen $f_\nu(x)$ und $\overline f_\nu(x)$ einander zugeordnet, so sind $\mathfrak N$ und $\overline{\mathfrak N}$ ähnlich-isomorph.}

Beweis durch Induktion: Ist $\sigma=1$, $n$ der Grad von $f_0(x)$ und $\overline f_0(x)$ und sind $j_0$, $\overline j_0$ die Hauptelemente, so erhält man alle Elemente von $\mathfrak N$ bzw. $\overline{\mathfrak N}$, indem man in allen Funktionen $g(x)$, $\overline g(x)$ aus $\K_0$ bzw. $\overline{\K}_0$, deren Grad kleiner als $n$ ist, $x$ durch $j_0$ bzw. $\overline j_0$ ersetzt. Den Funktionen $g(x)$ kommt in $\K_0$, den Funktionen $\overline g(x)$ in $\overline{\K}_0$ eine bestimmte Reihenfolge zu, und aus dem Umstande, daß $\pi_0$ eine ähnliche Beziehung zwischen $\K_0$ und $\overline{\K}_0$ darstellt, ergibt sich sofort, daß auch die durch $\pi_0$ vermittelte Beziehung zwischen den Funktionen aus $\K_0$ und $\overline{\K}_0$ eine ähnliche ist. Weiter folgt aus \S\,20, Def. 5, daß wir eine ähnliche $\pi_0$ umfassende Beziehung $\pi_1$ zwischen $\mathfrak N=\K_0(j_0)$ und $\overline{\mathfrak N}=\overline{\K}_0(\overline j_0)$ erhalten, wenn wir jedem Element $g(j_0)$ aus $\mathfrak N$ das Element $\overline g(\overline j_0)$ aus $\overline{\mathfrak N}$ entsprechen lassen, wobei $\overline g(x)$ die der Funktion $g(x)$ aus $\K_0$ in der Beziehung $\pi_0$ entsprechende Funktion aus $\overline{\K}_0$ bezeichnet. Da ferner $\pi_0$ einen Isomorphismus darstellt, so folgt, wie bei \S\,8, 1., daß auch die Beziehung $\pi_1$ zugleich isomorph ist. -- Ist $\sigma>1$ und bezeichnen wir mit $\K_\nu$, $\overline{\K}_\nu$ $(\nu<\sigma)$ die Hauptabschnitte von $\mathfrak N$ bzw. $\overline{\mathfrak N}$, so stellen $\K_\nu$ und $\overline{\K}_\nu$ normale und normal geordnete Erweiterungen von $\K_0$ und $\overline{\K}_0$ dar, deren Hauptfunktionen durch $\pi_0$ ineinander übergeführt werden, und wir dürfen als erwiesen ansehen, daß $\K_\nu$ und $\overline{\K}_\nu$ ähnlich-isomorph sind, d. h. daß die einzige zwischen $\K_\nu$ und $\overline{\K}_\nu$ bestehende ähnliche Beziehung $\pi_\nu$ zugleich isomorph ist. Ist $\mu<\nu<\sigma$, so sind $\K_\mu$ und $\overline{\K}_\mu$ ähnliche Abschnitte von $\K_\nu$ bzw. $\overline{\K}_\nu$, werden also durch $\pi_\nu$, ebenso wie durch $\pi_\mu$, ähnlich aufeinander abgebildet. Die Beziehung $\pi_\mu$ bildet also einen Teil von $\pi_\nu$. Daraus folgt, wenn $\sigma$ Limeszahl, also jedes Element von $\mathfrak N$ bzw. $\overline{\mathfrak N}$ in einem Körper $\K_\nu$ bzw. $\overline{\K}_\nu$ vorkommt, daß durch die Beziehungen $\pi_\nu$ eine ein-eindeutige Abbildung $\pi_\sigma$ von $\mathfrak N$ auf $\overline{\mathfrak N}$ hergestellt wird. Eine einfache Überlegung zeigt, daß die Abbildung $\pi_\sigma$ eine ähnliche ist, weil dies bei den Abbildungen $\pi_\nu$ der Fall ist; ebenso, daß $\pi_\sigma$ ein Isomorphismus ist, weil die Beziehungen $\pi_\nu$ Isomorphismen darstellen. -- Ist $\sigma$ keine Limeszahl, so hat man zunächst zwischen $\K_{\sigma-1}$ und $\overline{\K}_{\sigma-1}$ die ähnlich-isomorphe Beziehung $\pi_{\sigma-1}$, welche $\pi_0$ umfaßt, also die Hauptfunktionen $f_{\sigma-1}(x)$ und $\overline f_{\sigma-1}(x)$ einander zuordnet. Es wird ferner durch $\pi_{\sigma-1}$ der Hauptfaktor $\varphi_{\sigma-1}(x)$ als erster Faktor von $f_{\sigma-1}(x)$ innerhalb $\K_{\sigma-1}$ in den ersten Faktor von $\overline f_{\sigma-1}(x)$ innerhalb $\overline{\K}_{\sigma-1}$, d. i. $\overline\varphi_{\sigma-1}(x)$, verwandelt. Die Körper $\mathfrak N$ und $\overline{\mathfrak N}$ erweisen sich als normale und normal geordnete Erweiterungen von $\K_{\sigma-1}$ und $\overline{\K}_{\sigma-1}$ mit nur je einer Hauptfunktion $\varphi_{\sigma-1}(x)$ bzw. $\overline\varphi_{\sigma-1}(x)$. Da $\K_{\sigma-1}$ und $\overline{\K}_{\sigma-1}$ durch $\pi_{\sigma-1}$ ähnlich-isomorph aufeinander abgebildet werden und $\pi_{\sigma-1}$ zugleich die Hauptfunktionen ineinander überführt, so folgt das Bestehen der ähnlich-isomorphen Beziehung zwischen $\mathfrak N$ und $\overline{\mathfrak N}$ wie im Falle $\sigma=1$.

Lassen wir in Satz 5. $\K_0$ und $\overline{\K}_0$ zu einem und demselben wohlgeordneten Körper zusammenfallen, so ergibt sich:

\textbf{6.} \emph{Sind $\mathfrak N$ und $\overline{\mathfrak N}$ normale und normal geordnete Erweiterungen des wohlgeordneten Körpers $\K_0$ und gehört zu beiden dasselbe System $\{f_\nu(x)\}$ von Hauptfunktionen, so sind $\mathfrak N$ und $\overline{\mathfrak N}$ ähnlich-isomorph.}

\textbf{7.} \emph{Ist $\K_0$ ein wohlgeordneter Körper, $F=F_\sigma=\{f_\nu(x)\}$ $(0\leq\nu<\sigma)$ ein fortschreitendes System von Normalfunktionen aus $\K_0$, so gibt es normale und normal geordnete Erweiterungen von $\K_0$ mit dem Hauptfunktionensystem $F$.}

Beweis durch Induktion: Im Falle $\sigma=1$ können wir nach \S\,8, 1. eine Erweiterung $\mathfrak N$ bestimmen, welche zur Zerfällung von $f_0(x)$ gerade hinreicht. $\mathfrak N$ ist normal, und da auch $f_0(x)$ normal ist, so wird $\mathfrak N=\K_0(j_0)$, wenn $j_0$ irgendeine Nullstelle von $f_0(x)$ in $\mathfrak N$ ist. Ordnen wir $\mathfrak N$ in bezug auf $\K_0$ primitiv und so, daß $j_0$ Hauptelement wird, so ist sofort zu sehen, daß $\mathfrak N$ den Bedingungen des Satzes entspricht. -- Ist $\sigma>1$ und keine Limeszahl, so bezeichne $\K_{\sigma-1}$ eine normale und normal geordnete Erweiterung von $\K_0$ mit dem Hauptfunktionensystem $F_{\sigma-1}=\{f_\nu(x)\}$ $(0\leq\nu<\sigma-1)$. Von den irreduziblen Faktoren, welche die Normalfunktion $f_{\sigma-1}(x)$ im wohlgeordneten Körper $\K_{\sigma-1}$ besitzt, sei $\varphi_{\sigma-1}(x)$ die erste. Ihr Grad muß, da $\K_{\sigma-1}$ zur Zerfällung von $F_{\sigma-1}$ gerade hinreicht (Satz 1.) und das Normalfunktionensystem $F$ fortschreitet, größer als 1 sein. Erweitert man $\K_{\sigma-1}$ durch Adjunktion einer Nullstelle $j_{\sigma-1}$ von $\varphi_{\sigma-1}(x)$ zu einem Körper $\mathfrak N$ und ordnet man diesen primitiv in bezug auf $\K_{\sigma-1}$ mit $j_{\sigma-1}$ als Hauptelement, so entspricht $\mathfrak N$ offenbar den Bedingungen des Satzes 7. -- Ist $\sigma$ Limeszahl, so bezeichne $\K_\nu'$ $(0<\nu<\sigma)$ eine normale und normal geordnete Erweiterung von $\K_0$ mit dem Hauptfunktionensystem $F_\nu=\{f_\lambda(x)\}$ $(0\leq\lambda<\nu)$. Ist $\mu<\nu$, so besitzt $\K_\nu'$ einen Hauptabschnitt, der in bezug auf $\K_0$ normal und normal geordnet ist und das Hauptfunktionensystem $F_\mu$ hat. Nach 6. muß dieser Abschnitt zu $\K_\mu'$ ähnlich-isomorph sein. Somit ergibt sich: Von den Körpern der Folge $\K_0,\{\K_\nu'\}$ $(0<\nu<\sigma)$ ist jeder einem Abschnitt jedes späteren ähnlich-isomorph. Mithin gibt es (\S\,20, 4.) einen wohlgeordneten Körper $\overline{\mathfrak N}$ von der Beschaffenheit, daß $\K_0$ einem Abschnitt $\overline{\K}_0$ und jeder Körper $\K_\nu'$ $(0<\nu<\sigma)$ einem Abschnitt $\overline{\K}_\nu$ von $\overline{\mathfrak N}$ ähnlich-isomorph ist, und daß jedes Element aus $\overline{\mathfrak N}$ in Abschnitten $\overline{\K}_\nu$ vorkommt. Aus dem Umstande, daß zwischen $\K_0$ und $\overline{\K}_0$ eine ähnlich-isomorphe Beziehung besteht, folgt, daß man einen zu $\overline{\mathfrak N}$ ähnlich-isomorphen Körper $\mathfrak N$ erhält, der zugleich Erweiterungskörper von $\K_0$ ist, wenn man den Abschnitt $\overline{\K}_0$ von $\overline{\mathfrak N}$ durch das wohlgeordnete System $\K_0$ ersetzt und verfügt, daß jede Kompositionsgleichung $(a+b=c,\ a\cdot b=d)$ aus $\overline{\mathfrak N}$ zu einer Kompositionsgleichung aus $\mathfrak N$ werden soll, falls für die in ihr eventuell auftretenden Elemente von $\overline{\K}_0$ die entsprechenden von $\K_0$ substituiert werden. Durch diese Substitution geht jeder Abschnitt $\overline{\K}_\nu$ von $\overline{\mathfrak N}$ in einen zu ihm, also auch zu $\K_\nu'$, ähnlich-isomorphen Abschnitt $\K_\nu$ von $\mathfrak N$ über, und jedes Element von $\mathfrak N$ kommt in solchen Abschnitten $\K_\nu$ vor. Durch die ähnlich-isomorphe Beziehung $\pi_\nu$ zwischen $\K_\nu'$ und $\K_\nu$ werden alle Ordnungsbeziehungen und Gleichungen des einen Körpers in richtige Ordnungsbeziehungen und Gleichungen des andern übergeführt. Da ferner durch $\pi_\nu$ die Elemente von $\K_0$, also auch die Funktionen des Systems $F$ nicht alteriert werden, so stellt $\K_\nu$, wie $\K_\nu'$, eine normale und normal geordnete Erweiterung von $\K_0$ mit $F_\nu$ als Hauptfunktionensystem dar. Aus dem Umstande, daß der Körper $\mathfrak N$ sich aus den Elementen der Körperfolge $\{\K_\nu\}$ zusammensetzt, und daß jeder Körper dieser Folge Abschnitt aller späteren und Abschnitt von $\mathfrak N$ ist, ist aber sofort zu ersehen, daß auch $\mathfrak N$ selbst in bezug auf $\K_0$ normal geordnet ist, und daß $F$ das zugehörige System von Hauptfunktionen darstellt.

Aus dem bisherigen ergibt sich nun leicht der fundamentale Satz:

\textbf{8.} \emph{Ist $\K_0$ ein beliebiger Körper, $S$ ein beliebiges System ganzer Funktionen in $\K_0$, so gibt es eine und im wesentlichen nur eine Erweiterung von $\K_0$, welche gerade hinreichend groß ist, um alle Funktionen des Systems $S$ in Linearfaktoren zu zerfällen.}

\emph{Beweis:} Bei der hier behandelten Frage kann das System $S$ durch irgendein äquivalentes System ersetzt werden (\S\,18, 3.), und es darf daher angenommen werden, daß $S$ nur Normalfunktionen enthält (\S\,18, 6.). Ferner können wir uns das System $S$ wohlgeordnet denken und, falls es nicht fortschreitend ist, durch ein äquivalentes fortschreitendes Teilsystem ersetzen (\S\,20, 1.). Demnach dürfen wir von vornherein annehmen, daß ein fortschreitendes System von Normalfunktionen $F=\{f_\nu(x)\}$ vorliegt. Ferner ist es gestattet, den Körper $\K_0$ wohlgeordnet vorauszusetzen. Dann aber gibt es nach 7. eine normale und normal geordnete Erweiterung $\mathfrak N$ von $\K_0$ mit $F$ als Hauptfunktionensystem, und diese ist nach 1. zur vollständigen Zerfällung von $F$ gerade hinreichend. -- Nehmen wir nun an, daß noch eine zweite Erweiterung $\overline{\mathfrak N}$ von $\K_0$ vorliegt, welche zur vollständigen Zerfällung von $F$ gerade hinreicht. Dann haben wir die Äquivalenz der Erweiterungen $\mathfrak N$ und $\overline{\mathfrak N}$ nachzuweisen. Zu diesem Zwecke machen wir vom Auswahlprinzip Gebrauch. Wir denken uns nämlich aus den Nullstellen, welche die Funktionen $f_\nu(x)$ in $\overline{\mathfrak N}$ besitzen, je eine herausgegriffen, so daß wir ein System $\overline{\mathfrak J}=\{i_\nu\}$ haben, wo $i_\nu$ Nullstelle von $f_\nu(x)$ ist. Da es sich um Normalfunktionen handelt, wird $\overline{\mathfrak N}=\K_0(\overline{\mathfrak J})$. Nach 4. kann $\overline{\mathfrak N}$ in bezug auf $\K_0$ normal so geordnet werden, daß $F$ das System der Hauptfunktionen wird (und daß die Elemente $i_\nu$ ausgezeichnete Nullstellen der Hauptfunktionen werden. Dies letztere ist jetzt von keiner Bedeutung mehr, aber ohne diesen Teil konnte der Satz 4. nicht bewiesen werden). Somit sind $\mathfrak N$ und $\overline{\mathfrak N}$ normale und normal geordnete Erweiterungen von $\K_0$ mit demselben System von Hauptfunktionen, und es besteht zwischen ihnen eine ähnlich-isomorphe Beziehung $\pi$ (Satz 6.). Da durch diese jedes Element des gemeinsamen Abschnitts $\K_0$ sich selbst zugeordnet wird, so sind die Erweiterungen $\mathfrak N$ und $\overline{\mathfrak N}$ äquivalent.

Aus 8. folgt weiter durch Spezialisierung:

\textbf{9.} \emph{Jeder Körper läßt sich, und zwar im wesentlichen nur auf eine Art, algebraisch zu einem algebraisch abgeschlossenen Körper erweitern. -- Jedem beliebigen Körpertypus entspricht auf diese Weise ein bestimmter abgeschlossener.}

Verstehen wir nämlich unter $S$ das System aller ganzen Funktionen aus $\K_0$, so ist ein Erweiterungskörper von $\K_0$ dann und nur dann algebraisch abgeschlossen und in bezug auf $\K_0$ algebraisch, wenn er zur Zerfällung von $S$ gerade hinreicht. Hieraus folgt nach 8. der erste Teil von 9. -- Es sei nun irgendein Körpertypus gegeben, es seien $\K_0$ und $\K_0'$ zwei Körper dieses Typus und $\Omega$ und $\Omega'$ algebraische und algebraisch abgeschlossene Erweiterungen von $\K_0$ und $\K_0'$. Dann ist zu zeigen, daß $\Omega$ und $\Omega'$ von gleichem Typus sind. Zu diesem Zweck führen wir ein System $\mathfrak S$ von neuen Elementen ein, die wir den nicht zu $\K_0'$ gehörigen Elementen von $\Omega'$ eindeutig zuordnen. Diese Zuordnung gibt zusammen mit einem Isomorphismus $\pi_0$ zwischen $\K_0$ und $\K_0'$ eine Abbildung $\pi_1$ zwischen $\Omega'$ und $\overline\Omega=\K_0+\mathfrak S$. Verfügen wir, daß durch $\pi_1$ jede Kompositionsgleichung $a'+b'=c'$, $a'\cdot b'=d'$ aus $\Omega'$, auch wenn die Elemente $a'$, $b'$ nicht beide zu $\K_0'$ gehören, in eine Kompositionsgleichung aus $\overline\Omega$ übergehen soll, so wird $\overline\Omega$ ein zu $\Omega'$ isomorpher, mithin algebraisch abgeschlossener Körper und stellt eine algebraische Erweiterung von $\K_0$ dar. $\Omega$ und $\overline\Omega$ sind daher nach dem ersten bereits bewiesenen Teil von 9. äquivalente Erweiterungen von $\K_0$, und wir haben eine isomorphe Beziehung $\overline\pi_1$ zwischen $\Omega$ und $\overline\Omega$, welche die Elemente von $\K_0$ ungeändert läßt. Aus $\overline\pi_1$ und $\pi_1$ resultiert die nachzuweisende isomorphe Beziehung zwischen $\Omega$ und $\Omega'$. Dieselbe stellt sich als eine Erweiterung des zwischen $\K_0$ und $\K_0'$ bestehenden Isomorphismus $\pi_0$ dar.

\section*{\S\,22. Transzendente Erweiterungen. Algebraische Abhängigkeit zwischen Elementen eines Erweiterungskörpers.}

\emph{Jede nicht algebraische Erweiterung eines Körpers nennen wir transzendent.} -- Den folgenden Untersuchungen, welche uns zu einer übersichtlichen Einteilung aller Körper führen sollen, legen wir einen Körper $\K$ zugrunde und beziehen auf diesen Ausdrücke wie ,,Erweiterungskörper``, ,,transzendent`` u. s. f., sofern nichts anderes bemerkt ist. Ist $\mathfrak S$ ein System von Elementen eines Erweiterungskörpers $\Omega$, $a$ ein Element aus $\Omega$, so nennen wir $a$ \emph{algebraisch abhängig} von $\mathfrak S$ (in bezug auf $\K$), wenn $a$ in bezug auf den Körper $\K(\mathfrak S)$ algebraisch ist; ein Elementensystem $\mathfrak S'$ nennen wir algebraisch abhängig von $\mathfrak S$, wenn jedes Element aus $\mathfrak S'$ von $\mathfrak S$ algebraisch abhängt.

\textbf{1.} \emph{Hängt das Element $a$ vom System $\mathfrak S$ algebraisch ab, so gibt es ein endliches Teilsystem $\mathfrak S'$ von $\mathfrak S$, von welchem $a$ algebraisch abhängt.}

Denn aus unserer Voraussetzung folgt, daß $a$ Nullstelle einer ganzen Funktion $\varphi(x)$ des Körpers $\K(\mathfrak S)$ ist. Da $\varphi(x)$ nur endlich viele Koeffizienten hat, so gibt es ein endliches Teilsystem $\mathfrak S'$ von $\mathfrak S$ derart, daß diese Koeffizienten schon in $\K(\mathfrak S')$ vorkommen. Dann hängt aber $a$ von $\mathfrak S'$ algebraisch ab.

\textbf{2.} \emph{Hängt $\mathfrak S_3$ von $\mathfrak S_2$, $\mathfrak S_2$ von $\mathfrak S_1$ algebraisch ab, so ist $\mathfrak S_3$ algebraisch abhängig von $\mathfrak S_1$.}

Setzen wir nämlich $\K(\mathfrak S_1)=\K_1$, $\K(\mathfrak S_2)=\K_2$, so folgt aus unserer Voraussetzung, daß $\K_1(\mathfrak S_2)$ in bezug auf $\K_1$, $\K_2(\mathfrak S_3)$ in bezug auf $\K_2$ algebraisch ist (\S\,7, 7.). Da $\K_2$ Teilkörper von $\K_1(\mathfrak S_2)$ ist, so sind die Elemente von $\K_2(\mathfrak S_3)$ in bezug auf $\K_1(\mathfrak S_2)$, also auch in bezug auf $\K_1=\K(\mathfrak S_1)$ algebraisch (\S\,7, 9.). Mithin sind die Elemente von $\mathfrak S_3$ in bezug auf $\K(\mathfrak S_1)$ algebraisch, w. z. b. w.

Nennen wir daher zwei Systeme $\mathfrak S_1,\mathfrak S_2$ \emph{algebraisch äquivalent} (in bezug auf $\K$) oder kurz: äquivalent, wenn jedes vom andern algebraisch abhängt, so sind zwei Systeme, die einem dritten äquivalent sind, auch untereinander äquivalent u. s. f.

Wir nennen ein System $\mathfrak S$ von Elementen eines Erweiterungskörpers \emph{algebraisch reduzibel} (in bezug auf $\K$) oder kurz reduzibel, wenn es einem echten Teil von sich selbst äquivalent ist oder aus einem einzigen, in bezug auf $\K$ algebraischen Element besteht; andernfalls nennen wir $\mathfrak S$ (algebraisch) irreduzibel. Aus dieser Erklärung ist sofort zu sehen, daß in einem reduziblen System, welches mehr als ein Element enthält, wenigstens ein Element vorkommt, das von dem System der übrigen Elemente algebraisch abhängt. Hieraus ergibt sich in Verbindung mit 1. unmittelbar:

\textbf{3.} \emph{Jedes Teilsystem eines irreduziblen Systems ist irreduzibel.}

\textbf{4.} \emph{Jedes reduzible System enthält ein endliches reduzibles Teilsystem.}

\textbf{Rein transzendente Erweiterungen.} -- Eine Erweiterung eines Körpers nennen wir \emph{rein transzendent}, wenn sie durch Adjunktion eines irreduziblen Systems $\mathfrak S$ erhalten werden kann. Besteht $\mathfrak S$ nur aus endlich vielen Elementen $x_1,x_2,\ldots,x_n$, so haben wir es mit einer solchen Erweiterung zu tun, wie wir sie am Schluß von \S\,5 betrachtet haben: Der Körper $\K(\mathfrak S)$ besteht aus den rationalen Funktionen der Transzendenten oder Unbestimmten $x_1,x_2,\ldots,x_n$ mit Koeffizienten aus $\K$. Ist das System $\mathfrak S$ unendlich, so besteht $\K(\mathfrak S)$ aus den Elementen aller Körper $\K(\mathfrak S')$, worin $\mathfrak S'$ jedes endliche Teilsystem von $\mathfrak S$ bezeichnet. Die allgemeinste rein transzendente Erweiterung wird daher erhalten, indem man ein beliebiges System $\mathfrak S$ von unendlich vielen ,,Unbestimmten`` einführt. Jedes endliche Teilsystem $\mathfrak S'$ von $\mathfrak S$ liefert dann eine Erweiterung $\K(\mathfrak S')$ der früher betrachteten Art, und die Gesamtheit aller Körper $\K(\mathfrak S')$ konstituiert einen Körper (\S\,2, 2.), der aus $\K$ durch Adjunktion des irreduziblen Systems $\mathfrak S$ hervorgeht. -- Hat man zwei Erweiterungen $\K(\mathfrak S_1)$, $\K(\mathfrak S_2)$, sind $\mathfrak S_1$ und $\mathfrak S_2$ irreduzible Systeme, und ist die endliche oder unendlich große Anzahl ihrer Elemente (Kardinalzahl, Mächtigkeit) gleich, d. h. kann man zwischen ihren Elementen eine ein-eindeutige Beziehung $\varphi$ herstellen, so gibt es einen bestimmten Isomorphismus $\pi$ zwischen $\K(\mathfrak S_1)$ und $\K(\mathfrak S_2)$, durch welchen jedes Element von $\K$ sich selbst zugeordnet wird und die Elemente von $\mathfrak S_1$ und $\mathfrak S_2$ aufeinander in derselben Weise wie durch $\varphi$ bezogen werden. Dies wurde für den Fall endlicher Systeme in \S\,5 gezeigt, folgt aber daraus auf Grund einer einfachen Überlegung sofort für unendliche Systeme. Also:

\textbf{5.} \emph{Zwei rein transzendente Erweiterungen $\K(\mathfrak S_1)$, $\K(\mathfrak S_2)$ eines Körpers $\K$, bei denen die irreduziblen Systeme $\mathfrak S_1$, $\mathfrak S_2$ gleich viele Elemente enthalten, sind äquivalent.}

\textbf{6.} \emph{Wird ein irreduzibles System $\mathfrak S$ durch Hinzufügung eines Elementes $a$ reduzibel, so ist $a$ von $\mathfrak S$ algebraisch abhängig.}

Das reduzible System $\mathfrak S+a$ enthält nämlich ein endliches reduzibles System $\mathfrak T$. $a$ muß in $\mathfrak T$ vorkommen, weil sonst $\mathfrak T$ Teil von $\mathfrak S$, also irreduzibel wäre. Es sei $\mathfrak T=\mathfrak T'+a$, und es bestehe $\mathfrak T'$ aus den Elementen $x_1,x_2,\ldots,x_n$. Da $\mathfrak T'$ irreduzibel ist, so ist jeder Körper der Folge
\[
\K,
\quad \K_1=\K(x_1),
\quad \K_2=\K_1(x_2),
\quad \ldots,
\quad \K_n=\K_{n-1}(x_n)=\K(\mathfrak T')
\]
in bezug auf den vorangehenden transzendent. Wäre auch $\K_n(a)$ in bezug auf $\K_n$ transzendent, so wäre (\S\,5, 9.) jedes Element aus $\mathfrak T$ von dem System der übrigen Elemente algebraisch unabhängig, $\mathfrak T$ wäre irreduzibel. Da dies nicht der Fall ist, so ist $a$ in bezug auf $\K_n=\K(\mathfrak T')$ algebraisch, also von $\mathfrak T'$ und somit auch von $\mathfrak S$ algebraisch abhängig.

\textbf{7.} \emph{Ist $\mathfrak S$ ein (in bezug auf $\K$) irreduzibles System, das Element $a$ in bezug auf $\K$ transzendent, aber von $\mathfrak S$ algebraisch abhängig, so enthält $\mathfrak S$ ein bestimmtes endliches Teilsystem $\mathfrak T$ von folgender Beschaffenheit: $a$ ist von $\mathfrak T$ algebraisch abhängig; jedes Teilsystem von $\mathfrak S$, von welchem $a$ algebraisch abhängt, enthält das System $\mathfrak T$; wird irgendein Element aus $\mathfrak T$ durch $a$ ersetzt, so geht $\mathfrak S$ in ein äquivalentes irreduzibles System über; keinem der übrigen Elemente von $\mathfrak S$ kommt diese Eigenschaft zu.}

\emph{Beweis:} Es bezeichne $\mathfrak S'$ irgendein Teilsystem von $\mathfrak S$, von welchem $a$ algebraisch abhängt (das System $\mathfrak S$ selbst nicht ausgeschlossen). Dann besitzt $\mathfrak S'$ endliche Teilsysteme, von denen $a$ algebraisch abhängt. Es sei $\mathfrak T$ ein solches mit möglichst wenig Elementen. Die Elemente von $\mathfrak T$ seien $\alpha_1,\alpha_2,\ldots,\alpha_m$, und es werde mit $\mathfrak T_i$ $(i=1,\ldots,m)$ das System bezeichnet, welches nach Fortlassung von $\alpha_i$ aus $\mathfrak T$ zurückbleibt (im Falle $m=1$ das leere System). Da $a$ vom System $\mathfrak T$ algebraisch abhängig, von jedem der Systeme $\mathfrak T_i$ algebraisch unabhängig ist, so ist das System $\mathfrak T+a$ reduzibel, während die Systeme $\mathfrak T_i+a$ irreduzibel sind (Satz 6.). Dies gilt auch im Falle $m=1$, weil $a$ in bezug auf $\K$ transzendent ist. Da das irreduzible System $\mathfrak T_i+a$ durch Hinzufügung von $\alpha_i$ in das reduzible $\mathfrak T+a$ übergeht, so hängt $\alpha_i$ von $\mathfrak T_i+a$ algebraisch ab, und die Systeme $\mathfrak T_i+a$ und $\mathfrak T+a$ sind äquivalent. Ebenso sind $\mathfrak T$ und $\mathfrak T+a$ äquivalent. Somit sehen wir: \emph{Die $m+1$ Systeme $\mathfrak T_i+a$, $\mathfrak T$ und $\mathfrak T+a$, also auch einander, sind äquivalent.}

Wir wollen nun zeigen, daß jedes Teilsystem $\mathfrak S''$ von $\mathfrak S$, von dem $a$ algebraisch abhängt, das System $\mathfrak T$ enthalten muß. Zu diesem Zweck bezeichnen wir mit $\mathfrak U$ ein endliches aus möglichst wenig Elementen bestehendes Teilsystem von $\mathfrak S''$, von welchem $a$ algebraisch abhängt, und weisen die Identität von $\mathfrak U$ und $\mathfrak T$ nach. Es bestehe $\mathfrak U$ aus den Elementen $\beta_1,\ldots,\beta_n$, und es bezeichne $\mathfrak U_i$ $(i=1,\ldots,n)$ das nach Fortlassung von $\beta_i$ aus $\mathfrak U$ zurückbleibende System. Dann sind die Systeme $\mathfrak U_i+a$ und $\mathfrak U$ irreduzibel und einander sowie dem reduziblen System $\mathfrak U+a$ äquivalent. -- Wegen ihrer Eigenschaft als kleinste Systeme, von denen $a$ algebraisch abhängt, kann keins der Systeme $\mathfrak T$, $\mathfrak U$ ein echter Teil des andern sein. Wären sie also verschieden, so dürften wir annehmen, daß $\alpha_1$ nicht in $\mathfrak U$, $\beta_1$ nicht in $\mathfrak T$ vorkommt. Es sei nun $\mathfrak B$ dasjenige Teilsystem von $\mathfrak S$, welches alle Elemente von $\mathfrak T$ und alle von $\mathfrak U$, aber sonst kein Element enthält. Da $\mathfrak T_1+a$ mit dem in $\mathfrak B$ enthaltenen System $\mathfrak T$ äquivalent ist, so erhält man aus $\mathfrak B$ ein zu $\mathfrak B$ äquivalentes System $\mathfrak B'$, wenn man $\alpha_1$ durch $a$ ersetzt. $\mathfrak B'$ enthält das System $\mathfrak U+a$, und da dieses $\mathfrak U$ äquivalent ist, so erhält man ein $\mathfrak B'$ äquivalentes System $\mathfrak B''$, wenn man aus $\mathfrak B'$ das Element $a$ fortläßt. Hiernach würde $\mathfrak B$ einem echten Teil $\mathfrak B''$ von sich äquivalent, also reduzibel sein. Das ist aber nicht möglich, weil $\mathfrak B$ Teil des irreduziblen Systems $\mathfrak S$ ist.

\emph{Es muß also in der Tat $\mathfrak T$ in jedem Teilsystem von $\mathfrak S$ enthalten sein, von welchem $a$ algebraisch abhängt.} Bezeichnet daher $\mathfrak S_i$ $(i=1,\ldots,m)$ das System, welches nach Fortlassung von $\alpha_i$ aus $\mathfrak S$ zurückbleibt, so ist $a$ von dem irreduziblen System $\mathfrak S_i$ algebraisch unabhängig, das System $\mathfrak S_i+a$ also ebenfalls irreduzibel. $\mathfrak S_i+a$ geht aber aus $\mathfrak S$ hervor, indem man $\alpha_i$ durch $a$ oder, was auf dasselbe hinauskommt, $\mathfrak T=\mathfrak T_i+\alpha_i$ durch das äquivalente System $\mathfrak T_i+a$ ersetzt. Es sind daher $\mathfrak S$ und $\mathfrak S_i+a$ äquivalente Systeme. Also: \emph{Man erhält ein irreduzibles und $\mathfrak S$ äquivalentes System, wenn man eins der Elemente von $\mathfrak T$ durch $a$ ersetzt.} -- Ist hingegen $\gamma$ ein nicht in $\mathfrak T$ vorkommendes Element von $\mathfrak S$, nach dessen Fortlassung $\mathfrak S'$ zurückbleibt, so enthält $\mathfrak S'$ das System $\mathfrak T$, $a$ hängt also von $\mathfrak S'$ algebraisch ab, und das System $\mathfrak S'+a$, welches aus $\mathfrak S$ hervorgeht, wenn $\gamma$ durch $a$ ersetzt wird, ist also reduzibel. Es ist dem System $\mathfrak S'$ äquivalent, und da dieses als echter Teil des irreduziblen Systems $\mathfrak S$ nicht $\mathfrak S$ äquivalent sein kann, so ist auch $\mathfrak S'+a$ nicht $\mathfrak S$ äquivalent. Damit ist der Beweis von 7. in allen Teilen erbracht.

\textbf{8.} \emph{Es seien $\mathfrak U$ und $\mathfrak B$ endliche irreduzible Systeme von $m$ bzw. $n$ Elementen; es sei $n\leq m$ und $\mathfrak B$ algebraisch abhängig von $\mathfrak U$. Dann sind im Falle $m=n$ die Systeme $\mathfrak U$ und $\mathfrak B$ äquivalent, im Falle $n<m$ aber ist $\mathfrak U$ einem irreduziblen System äquivalent, welches aus $\mathfrak B$ und $m-n$ Elementen aus $\mathfrak U$ besteht.}

Es seien nämlich $u_1,\ldots,u_m$ die Elemente von $\mathfrak U$, $v_1,\ldots,v_n$ die von $\mathfrak B$. Da $v_1$ in bezug auf $\K$ transzendent ist und von $\mathfrak U$ algebraisch abhängt, so gibt es nach 7. in $\mathfrak U$ wenigstens ein Element, nach dessen Ersatz durch $v_1$ aus $\mathfrak U$ ein zu $\mathfrak U$ äquivalentes, irreduzibles System hervorgeht. Damit ist die Behauptung für den Fall $n=1$ erwiesen. Für $n>1$ ist der Beweis durch die vollständige Induktion zu führen. Da das System $\mathfrak B'$, welches aus den Elementen $v_1,\ldots,v_{n-1}$ besteht, irreduzibel ist und von $\mathfrak U$ algebraisch abhängt, so ist als bewiesen zu erachten, daß $\mathfrak U$ einem irreduziblen System $\mathfrak U'$ äquivalent ist, welches aus $\mathfrak B'$ und $m-n+1$ Elementen aus $\mathfrak U$ besteht. $v_n$ hängt von $\mathfrak U'$, nicht aber von $\mathfrak B'$ algebraisch ab. Ist daher $\mathfrak T$ das kleinste Teilsystem von $\mathfrak U'$, von welchem $v_n$ algebraisch abhängt, so enthält $\mathfrak T$ wenigstens eines der $m-n+1$ Elemente aus $\mathfrak U$. Ersetzen wir ein solches durch $v_n$, so wird aus $\mathfrak U'$ ein irreduzibles, $\mathfrak U$ äquivalentes System, welches aus $\mathfrak B$ und $m-n$ Elementen von $\mathfrak U$ besteht. Im Falle $m=n$ sind also $\mathfrak U$ und $\mathfrak B$ äquivalent.

Wir ersehen daraus auch, daß ein irreduzibles, von $\mathfrak U$ algebraisch abhängiges System $\mathfrak B$ nicht mehr als $m$ Elemente enthalten kann. Andernfalls würde ein echter aus $m$ Elementen bestehender Teil $\mathfrak B'$ von $\mathfrak B$ nach dem eben Bewiesenen mit $\mathfrak U$ äquivalent sein. Es wäre $\mathfrak B$ von $\mathfrak B'$ algebraisch abhängig, also reduzibel. Also:

\textbf{9.} \emph{Ein irreduzibles System $\mathfrak B$, welches von einem endlichen System $\mathfrak U$ algebraisch abhängig ist, kann nicht mehr Elemente enthalten als dieses; im Falle gleicher Elementenzahl sind beide Systeme äquivalent.}
